(1) 界面平行方向波数守恒，\(k_x=n_1k_0\sin\theta\)。高折射率棱镜中 \(k_z=\sqrt{n_1^2k_0^2-k_x^2}=n_1k_0\cos\theta>0\)，低折射率层中 \(k_{z,2}=i\kappa\)，\(\kappa=\sqrt{k_x^2-n_2^2k_0^2}>0\)，故间隙内一般解是 \(Ae^{-\kappa z}+Be^{+\kappa z}\) 而非只能保留单向衰减项。
(2) 可将TE Helmholtz方程的分段法向问题等价成一维对称矩形势垒匹配。分别施加两界面的场与导数连续，得到左右介质相同的精确功率透射系数 \(T=\left[1+\frac{(k_z^2+\kappa^2)^2}{4k_z^2\kappa^2}\sinh^2(\kappa d)\right]^{-1}\)。系数非负，故 \(0<T\le1\)，无吸收时反射率 \(R=1-T\)。
(3) 当 \(\kappa d\gg1\)，有 \(\sinh^2(\kappa d)\sim e^{2\kappa d}/4\)，于是 \(T\sim[16k_z^2\kappa^2/(k_z^2+\kappa^2)^2]e^{-2\kappa d}\)。当 \(d\to0\) 时 \(\sinh(\kappa d)\to0\)，精确式给 \(T\to1\)，因为两棱镜实际上接触成为相同介质。
(4) 若低折射率区域延伸到 \(+\infty\)，有界条件迫使倏逝解仅为 \(e^{-\kappa z}\)，其 \(E^*\partial_zE\) 为纯实数，平均法向通量为零。有限间隙中右侧出射传播波要求同时存在衰减和增长的倏逝振幅分量，它们的相对复相位形成非零法向能流；这是边界耦合造成的传输，不违反无源介质的能量守恒。
